\section{相关工作}
\label{sec:related}

\subsection{医学图像翻译}
非配对跨模态合成经历了从对抗式翻译的演进——CycleGAN~\cite{zhu2017cyclegan} 与
pix2pix~\cite{isola2017pix2pix} 确立了非配对与配对两种范式，随后是多对比度 MRI 专精模型，
如条件 GAN~\cite{dar2019cgan} 与 CNN--Transformer 混合架构
（ResViT~\cite{dalmaz2022resvit}）——再到扩散类模型。SynDiff~\cite{ozbey2023syndiff}
以循环一致性先验的对抗式扩散实现无监督医学翻译；I2I-Galip~\cite{korkmaz2025i2igalip}
用视觉--语言模型替换判别器以提升效率。平行地，扩散桥不经过各向同性噪声瓶颈，
而是直接连接源与目标两个图像分布：DDPM~\cite{ho2020ddpm} 及其基于分数的推广~\cite{song2021scoresde}
提供了数学骨架，谱系中包括基于分数的 NCSN~\cite{song2019ncsn}、改进的调度~\cite{nichol2021iddpm}、
图像到图像薛定谔桥~\cite{liu2023i2sb}、带数据一致性的直接扩散桥~\cite{chung2023directbridge}、
对偶隐式桥~\cite{su2023ddib}、神经薛定谔桥~\cite{kim2024nsb}、直接迭代反演~\cite{delbracio2023indi}
以及三维桥理论~\cite{zhang2025bridge3d}。SelfRDB~\cite{arslan2025selfrdb} 将软先验桥与
自洽递归采样结合，报告了最先进性能。上述工作均以翻译质量为导向，不提供输出可信度信息。
\textbf{与之相对，本文在不改动翻译模型的前提下，从其推理过程内部提取可信度证据。}

\subsection{医学影像不确定性量化}
模型级方法包括 MC-Dropout~\cite{gal2016mcdropout}、深度集成~\cite{lakshminarayanan2017ensembles}
与测试时增强~\cite{ayhan2018tta}；不确定性可分为偶然性与认知性两类~\cite{kendall2017uncertainties}。
扩散模型语境下，集成式不确定性已用于分割~\cite{wolleb2022diffusionensembles} 与
显微图像复原~\cite{osuna2024microscopy}：对同一输入采样多条链，以链间分歧度量困难区域。
本文继承该思路但有三点不同：(i) 我们的集成成员来自扩散桥自身原生的随机终点与隐变量，
无需注入额外随机性；(ii) 我们配对了一个与采样方差正交的免费信号——CReF，
它只存在于自洽递归采样器中；(iii) 我们将信号与临床 QC 场景对接，但图像级聚合被证明不可靠（\S4.6），本文定位为像素级排序工具。
值得注意的是，部分桥模型将这种随机性当作缺陷主动抑制~\cite{song2024ccbm}；
本文论证：在可信度评估的任务里，它是资产。

\subsection{生成影像的幻觉评估与可信放射 AI}
与本文目标最接近的 Hallucination Index~\cite{tivnan2024hallucinationindex}
将生成式重建模型的编造内容量化为对重复采样计算的\emph{图像级}指标；
MUMINS~\cite{oliveras2026mumins} 以元数据为条件进行医学图像合成并训练不确定性感知
预测头输出空间不确定性图；\cite{jin2026trustworthy} 的测量几何分析证明了测量算子
何时能够区分不同解，从算子设计视角解释幻觉。这些方法或需要图像粒度的重复采样，
或需要改动架构并重训，或需要事前的算子设计。TrustBridge 则在\emph{现有采样器内部}
提取像素级证据，不需要架构改动或标注，并进一步支持规则化门控：我们的框架与
RAISE~\cite{cardoso2023raise} 这类输出护栏生命周期互补，为
AI-RADS~\cite{russe2026airads} 这类报告框架所呼唤的客观像素级证据提供供给，
同时保持全自动。

\paragraph{超越采样分歧：语义与病理感知信号。}
分歧式证据对采样器\emph{持续一致地犯}的错误是盲的——我们在 \S4 量化了这一局限。
三条并行研究线从评估侧攻击该问题，并构成我们这类信任层的自然下一步。
CRAFT~\cite{chung2026craft} 用基础模型的临床对齐分数（CAS）按临床标准为生成图像打分，
抑制幻觉样内容；ViCTr~\cite{susladkar2025victr} 通过 EWC 正则的 LoRA 适配实现
分级严重度控制的病理感知合成；DDWM~\cite{li2025ddwm} 只学习目标域分布、
经小波分解融合域不变解剖，实现零样本翻译。三者或依赖基础模型、或需重训、
或需架构改动；它们与采样器原生证据互补，我们在 \S6 展望中讨论。

\paragraph{方法对照。}
\begin{table}[h]
\centering\small
\caption{不确定性/幻觉评估方法对照（"额外前向"指相对单次基础推理的生成器前传次数）。
TrustBridge 一行的"无额外前传"限定于当前 SelfRDB 实现——其递归结构为常规 DDPM 所无。}
\label{tab:rw}
\begin{tabular}{lcccc}
\toprule
方法 & 需要训练 & 额外模型 & 额外前向 & 输出粒度 \\
\midrule
MC-Dropout~\cite{gal2016mcdropout} & 需（启用 dropout 重训） & 否 & 多次 & 像素/图像 \\
深度集成~\cite{lakshminarayanan2017ensembles} & 需（多模型） & 是 & 多次 & 像素/图像 \\
TTA~\cite{ayhan2018tta} & 否 & 否 & 多次 & 像素/图像 \\
扩散集成~\cite{wolleb2022diffusionensembles,osuna2024microscopy} & 否（复用模型） & 否 & 多次 & 像素/图像 \\
Hallucination Index~\cite{tivnan2024hallucinationindex} & 视实现 & 视实现 & 多次 & 图像级 \\
TrustBridge（本文） & 否 & 否 & \textbf{无额外}（单链） & 像素级 \\
\bottomrule
\end{tabular}
\end{table}
